
\documentclass[aspectratio=169,11pt]{beamer}

\usetheme{Madrid}
\usecolortheme{default}
\setbeamertemplate{navigation symbols}{}
\setbeamertemplate{footline}[frame number]

\usepackage[T1]{fontenc}
\usepackage[utf8]{inputenc}
\usepackage[french]{babel}
\usepackage{lmodern}
\usepackage{tikz}
\usetikzlibrary{arrows.meta,positioning,fit,calc,shapes.geometric}
\usepackage{booktabs}
\usepackage{tabularx}
\usepackage{array}
\usepackage{ragged2e}
\usepackage{hyperref}

\newcommand{\method}[1]{\textbf{#1}}
\newcommand{\source}[1]{\scriptsize\textcolor{gray}{Source : #1}}

\title{De la maquette à une représentation hiérarchique}
\subtitle{Étude comparative de quatre approches de génération de code frontend}

\author{KOUONGA TADIE KENNEDY}
\institute{Universite de DOUALA}
\date{Jeudi 6 août 2026}
\begin{document}

% ------------------------------------------------
\begin{frame}
  \titlepage
\end{frame}

% ------------------------------------------------
\begin{frame}{Question de recherche}
\begin{block}{Problème étudié}
Comment passer d'une représentation visuelle d'une interface
(maquette, capture d'écran, croquis, etc.) à une
\textbf{représentation structurée et hiérarchique} qui pourra ensuite
servir à générer automatiquement du code frontend ?
\end{block}

\vspace{0.5cm}

\centering
\begin{tikzpicture}[node distance=1.1cm,align=center, every node/.style={font=\small}]
\node[draw, rounded corners, minimum width=2.7cm, minimum height=1cm] (input)
{Maquette /\\capture / croquis};
\node[draw, rounded corners, minimum width=3.0cm, minimum height=1cm, right=of input]
(model) {Modèle\\hiérarchique};
\node[draw, rounded corners, minimum width=2.7cm, minimum height=1cm, right=of model]
(code) {Code frontend};
\draw[-{Latex},thick] (input) -- (model);
\draw[-{Latex},thick] (model) -- (code);
\end{tikzpicture}

\vspace{0.3cm}
\source{Problématique de recherche ; travaux étudiés}
\end{frame}

% ------------------------------------------------
\begin{frame}{Pourquoi passer par une représentation intermédiaire ?}
\begin{itemize}
  \item Une image décrit surtout \textbf{ce que l'on voit}.
  \item Le code doit également exprimer \textbf{comment les éléments sont organisés}.
  \item Il faut donc retrouver :
  \begin{itemize}
    \item les éléments de l'interface ;
    \item leurs relations ;
    \item leur organisation spatiale ;
    \item et, selon la méthode, leurs comportements de mise en page.
  \end{itemize}
  \item Les quatre travaux étudiés ne construisent cependant pas tous le même type de représentation.
\end{itemize}

\vspace{0.25cm}
\begin{alertblock}{Point important}
Dans cette étude, « représentation hiérarchique » est utilisé au sens large :
arbre de composants, arbre de mise en page, arbre de subdivisions ou
structure de mise en page avec contraintes.
\end{alertblock}
\end{frame}

% ------------------------------------------------
\begin{frame}{Les quatre approches étudiées}
\small
\begin{tabularx}{\textwidth}{>{\raggedright\arraybackslash}p{3.0cm} X}
\toprule
\textbf{Approche} & \textbf{Idée principale pour structurer l'interface}\\
\midrule
Prototype2Code & Reconstruire une hiérarchie à partir des relations de contenance, puis organiser les éléments en lignes/colonnes.\\
DCGen & Diviser récursivement l'image en segments ; la division elle-même forme un arbre.\\
ScreenCoder & Détecter et nommer les grandes régions, puis construire un arbre de mise en page à partir de règles de conception web.\\
DesignCoder & Diviser l'interface, extraire le sens des éléments, puis les regrouper dans un arbre de composants.\\
\bottomrule
\end{tabularx}

\vspace{0.25cm}
\source{Synthèse des quatre documents fournis}
\end{frame}

% =================================================
\section{Prototype2Code}

\begin{frame}{1. Prototype2Code : partir des relations de contenance}
\begin{block}{Entrée}
Une maquette contenant des éléments visuels et des informations de conception.
\end{block}

\begin{enumerate}
  \item Nettoyer les informations issues de la maquette.
  \item Identifier les groupes et les éléments terminaux.
  \item Déterminer quels éléments sont contenus dans quels groupes.
  \item Construire un arbre de groupes imbriqués.
  \item Organiser ensuite les enfants en \textbf{Row} ou \textbf{Column}.
\end{enumerate}

\source{Document fourni sur Prototype2Code}
\end{frame}

% ------------------------------------------------
\begin{frame}{Prototype2Code : construction du modèle}
\centering
\begin{tikzpicture}[
  node/.style={draw, rounded corners, minimum height=0.65cm, minimum width=1.8cm,align=center, font=\scriptsize},
  arrow/.style={-{Latex}, thick}
]
\node[node] (groups) {Groupes détectés};
\node[node, right=1.5cm of groups] (contain) {Relations de\\contenance};
\node[node, right=1.5cm of contain] (tree) {Arbre initial};
\node[node, below=1.0cm of tree] (rowcol) {Segmentation\\Row / Column};
\node[node, left=1.5cm of rowcol] (responsive) {Mise en page\\adaptative};

\draw[arrow] (groups) -- (contain);
\draw[arrow] (contain) -- (tree);
\draw[arrow] (tree) -- (rowcol);
\draw[arrow] (rowcol) -- (responsive);
\end{tikzpicture}

\vspace{0.5cm}
\begin{itemize}
  \item La hiérarchie repose d'abord sur la \textbf{géométrie de la maquette}.
  \item Les éléments sont placés sous le plus petit groupe qui les contient.
  \item Les chevauchements peuvent conduire à un nœud utilisant un positionnement absolu.
  \item La seconde phase cherche une organisation compatible avec Flexbox.
\end{itemize}

\source{Document Prototype2Code fourni, section sur l'UI Layout Tree}
\end{frame}

% =================================================
\section{DCGen}

\begin{frame}{2. DCGen : construire un arbre en divisant l'image}
\begin{block}{Idée centrale}
Au lieu d'essayer de comprendre toute la page en une seule fois,
DCGen cherche des \textbf{lignes de séparation} et découpe récursivement
la capture en parties plus simples.
\end{block}

\begin{enumerate}
  \item Détecter des séparations horizontales et verticales.
  \item Découper l'image en segments.
  \item Répéter le découpage sur les segments obtenus.
  \item Stocker ces divisions sous forme d'un arbre.
  \item Générer le code des petites parties puis les réassembler.
\end{enumerate}

\source{Wan et al., 2025}
\end{frame}

% ------------------------------------------------
\begin{frame}{DCGen : la hiérarchie représente les divisions}
\centering
\begin{tikzpicture}[
  every node/.style={draw, rounded corners, minimum width=1.8cm, minimum height=0.55cm, font=\scriptsize, align=center},
  edge from parent/.style={draw,-{Latex},thick},
  level 1/.style={sibling distance=4.0cm},
  level 2/.style={sibling distance=2.0cm}
]
\node {Capture complète}
 child {node {Segment A}
   child {node {A1}}
   child {node {A2}}
 }
 child {node {Segment B}
   child {node {B1}}
   child {node {B2}}
 };
\end{tikzpicture}

\vspace{0.35cm}
\begin{itemize}
  \item Le nœud parent représente une zone avant sa subdivision.
  \item Les feuilles représentent les plus petites zones obtenues.
  \item L'arbre sert ensuite à reconstruire progressivement le code de la page.
\end{itemize}

\begin{alertblock}{Différence importante}
Ici, l'arbre décrit d'abord une \textbf{décomposition spatiale de l'image},
et non une compréhension sémantique complète des composants.
\end{alertblock}
\end{frame}



% =================================================
\section{ScreenCoder}

\begin{frame}{3. ScreenCoder : détecter, planifier}

\begin{block}{Idée centrale}
ScreenCoder décompose la génération d'interface en trois étapes :
\textbf{comprendre visuellement}, \textbf{construire la structure}, puis
\textbf{générer le code}.
\end{block}

\begin{enumerate}
\item \textbf{Grounding : comprendre l'interface avec un MLLM}


Le \textbf{MLLM} analyse la capture d'écran pour identifier les
regroupements visuels importants et leur attribuer un rôle :
\textit{header, sidebar, navigation, contenu, etc.}

Pour chaque région détectée, il fournit une
\textbf{bounding box} définie par :
\[
(x,y,width,height)
\]

\item \textbf{Planning : construire l'arbre de mise en page}

Les régions détectées sont ensuite organisées à partir de leurs
\textbf{positions et dimensions}. Les relations spatiales entre les
bounding boxes permettent de déterminer la hiérarchie :
quelles régions contiennent d'autres régions et quelles régions sont
placées côte à côte.

On obtient ainsi un \textbf{arbre de mise en page hiérarchique},
proche de la structure d'un DOM.



\end{enumerate}

\begin{center}
\textit{Capture d'écran}
$\rightarrow$
\textbf{MLLM + bounding boxes}
$\rightarrow$
\textbf{arbre de mise en page}
$\rightarrow$
\textbf{HTML/CSS}
\end{center}

\source{Jiang et al., ScreenCoder}
\end{frame}


% ------------------------------------------------
\begin{frame}{ScreenCoder : comment le modèle est construit ?}
\centering
\begin{tikzpicture}[
  node/.style={draw, rounded corners, minimum width=2.6cm, minimum height=0.75cm, font=\scriptsize, align=center},
  arrow/.style={-{Latex}, thick}
]
\node[node] (input) {Screenshot /\\Design sketch};
\node[node, right=1.0cm of input] (ground) {Régions +\\labels + positions};
\node[node, right=1.0cm of ground] (plan) {Arbre de\\mise en page};
\node[node, right=1.0cm of plan] (code) {HTML + CSS};

\draw[arrow] (input) -- (ground);
\draw[arrow] (ground) -- (plan);
\draw[arrow] (plan) -- (code);
\end{tikzpicture}

\vspace{0.45cm}
\begin{itemize}
  \item Le planning convertit ces informations 2D en une structure de type DOM.
  \item Les coordonnées sont normalisées en pourcentage pour aider à préserver la structure.
  \item Les conteneurs internes peuvent être organisés avec des mécanismes de mise en page comme Grid.
\end{itemize}

\source{ScreenCoder, sections Grounding et Planning}
\end{frame}

% =================================================
\section{DesignCoder}
\begin{frame}{4. DesignCoder : comprendre les composants avant l'arbre}

\begin{block}{Idée centrale}
DesignCoder ne se limite pas à la position des éléments :
il cherche à comprendre leur \textbf{rôle} et leur \textbf{contexte}
avant de construire la hiérarchie des composants.
\end{block}

\begin{enumerate}

\item \textbf{Entrée : image + métadonnées}

La maquette fournit une image ainsi que les métadonnées de chaque
élément :
[
(id,\ type,\ bbox,\ style)
]

La \textit{bbox} indique notamment la position et les dimensions
de l'élément.

\item \textbf{UI Division : découper l'interface en régions}

Le \textbf{MLLM analyse l'image et les métadonnées} pour diviser
l'interface en sous-régions cohérentes.

Une région correspond à un \textbf{ensemble d'éléments liés par leur
rôle ou leur fonction}, par exemple une barre de recherche,
une bannière ou une barre de navigation.

\item \textbf{Semantic Extraction : comprendre chaque élément}

Le MLLM observe chaque élément dans son contexte visuel et décrit
sa \textbf{fonction et son rôle}.

\begin{center}
\textit{Exemple : « cette icône représente une recherche »}
\end{center}

\end{enumerate}

\vspace{0.2cm}

\begin{center}
\textbf{Image + métadonnées}
$\rightarrow$
\textbf{Régions}
$\rightarrow$
\textbf{Sémantique des éléments}
\end{center}

\source{Chen et al., DesignCoder}
\end{frame}


\begin{frame}{DesignCoder : du regroupement à l'arbre de composants}

\begin{block}{Component Grouping}
Après avoir identifié les régions et compris le rôle des éléments,
le MLLM regroupe les éléments selon plusieurs relations.
\end{block}

\begin{itemize}

\item \textbf{Relations spatiales} :
position, proximité, alignement et organisation des éléments.

\item \textbf{Relations sémantiques} :
éléments ayant un rôle ou un sens commun.

\item \textbf{Relations fonctionnelles} :
éléments participant à une même fonctionnalité.

\end{itemize}

\vspace{0.15cm}

\begin{block}{Construction de l'arbre}
Le regroupement transforme la liste plate des éléments en une
\textbf{hiérarchie de composants}.

Les composants regroupés deviennent des \textbf{nœuds parents},
et les éléments qui les composent deviennent leurs \textbf{enfants}.
\end{block}

\vspace{0.15cm}


\vspace{0.1cm}

\begin{center}
\textbf{Position / bbox}
$+$
\textbf{contexte visuel}
$+$
\textbf{sémantique}
$+$
\textbf{fonction}
$\Rightarrow$
\textbf{regroupement}
$\Rightarrow$
\textbf{arbre}
\end{center}

\source{Chen et al., DesignCoder}
\end{frame}


% ------------------------------------------------

\begin{frame}{DesignCoder : comment les éléments deviennent un arbre ?}

\centering

\begin{tikzpicture}[
node distance=0.65cm and 0.55cm,
box/.style={
draw,
rounded corners,
minimum width=2.45cm,
minimum height=0.72cm,
font=\scriptsize,
align=center
},
arrow/.style={
-{Latex},
thick
}
]

\node[box] (input) {Image + métadonnées};

\node[box, below=of input] (division) {

Régions sémantiques
};

\node[box, right=1.0cm of division] (semantic) {

Rôle de chaque élément
};

\node[box, right=1.0cm of semantic] (group) {
 Relations entre éléments
};

\node[box, below=of group] (tree) {
Hiérarchie des composants
};

\draw[arrow] (input) -- (division);
\draw[arrow] (division) -- (semantic);
\draw[arrow] (semantic) -- (group);
\draw[arrow] (group) -- (tree);

\end{tikzpicture}

\vspace{0.35cm}

\begin{block}{Ce qui détermine le regroupement}
\centering
\textbf{Position / bbox}
\quad + \quad
\textbf{Sens visuel}
\quad + \quad
\textbf{Fonction}
\quad
$\Longrightarrow$
\quad
\textbf{Composant}
\end{block}

\begin{itemize}
\item La \textbf{bbox} fournit la position et les dimensions des éléments.
\item Le \textbf{MLLM} interprète leur rôle et leur contexte visuel.
\item Les relations \textbf{spatiales, sémantiques et fonctionnelles}
guident le regroupement.
\item Le regroupement produit finalement un
\textbf{arbre hiérarchique de composants}.
\end{itemize}

\source{DesignCoder, UI Grouping Chain}
\end{frame}




% =================================================
\section{Comparaison}

\begin{frame}{Comparaison : quelle information permet de construire la hiérarchie ?}
\scriptsize
\begin{tabularx}{\textwidth}{>{\raggedright\arraybackslash}p{2.6cm} X X}
\toprule
\textbf{Méthode} & \textbf{Principe de construction} & \textbf{Information dominante}\\
\midrule
Prototype2Code &
Relations de contenance puis Row/Column &
Géométrie + structure de la maquette\\
DCGen &
Découpage récursif de l'image &
Séparations et organisation spatiale\\
ScreenCoder &
Détection de régions puis mapping vers un arbre &
Régions nommées + positions + règles frontend\\
DesignCoder &
Division + sémantique + regroupement &
Sémantique + relations spatiales/fonctionnelles + métadonnées\\
\bottomrule
\end{tabularx}

\vspace{0.35cm}
\begin{block}{Constat}
Les méthodes cherchent toutes à retrouver une structure, mais
\textbf{elles ne cherchent pas exactement la même structure}.
\end{block}
\end{frame}

% ------------------------------------------------
\begin{frame}{Comparaison : nature de la représentation obtenue}
\scriptsize
\begin{tabularx}{\textwidth}{>{\raggedright\arraybackslash}p{2.6cm} X X}
\toprule
\textbf{Méthode} & \textbf{Représentation intermédiaire} & \textbf{Rôle pour le code}\\
\midrule
Prototype2Code &
Arbre de mise en page avec groupes et Row/Column &
Guide l'imbrication HTML et les styles Flexbox/positionnement\\
DCGen &
Arbre de subdivisions de l'image &
Guide l'assemblage progressif du code\\
ScreenCoder &
Arbre de layout de type DOM &
arbre de composants pour HTML/CSS\\
DesignCoder &
Arbre hiérarchique de composants enrichi par les styles &
Guide la génération structurée des composants et du CSS\\
\bottomrule
\end{tabularx}

\vspace{0.3cm}
\source{Synthèse des documents étudiés}
\end{frame}

% ------------------------------------------------
\begin{frame}{Comparaison : points forts et limites}
\scriptsize
\begin{tabularx}{\textwidth}{>{\raggedright\arraybackslash}p{2.4cm} X X}
\toprule
\textbf{Méthode} & \textbf{Point fort pour notre problème} & \textbf{Limite / difficulté}\\
\midrule
Prototype2Code &
Modèle directement orienté vers la mise en page et le code &
Dépend fortement des informations de conception et de règles géométriques\\
DCGen &
Obtient une hiérarchie même à partir d'une capture d'écran &
L'arbre représente surtout les divisions spatiales, pas nécessairement les composants sémantiques\\
ScreenCoder &
Sépare clairement perception, structure et génération &
Le planning repose sur des connaissances de mise en page web\\
DesignCoder &
Combine géométrie, sémantique et hiérarchie de composants &
Dépend d'un traitement multimodal et d'informations de design riches\\
\bottomrule
\end{tabularx}
\end{frame}

% ------------------------------------------------
\begin{frame}{Ce que les quatre méthodes ont en commun}
\begin{enumerate}
  \item \textbf{Identifier des unités de l'interface} : éléments, régions ou widgets.
  \item \textbf{Déterminer des relations} entre ces unités.
  \item \textbf{Construire une structure} permettant de ne plus traiter l'interface comme une simple image plate.
  \item \textbf{Utiliser cette structure} pour faciliter la génération ou la reconstruction du code.
\end{enumerate}

\vspace{0.35cm}
\centering
\begin{tikzpicture}[node distance=0.9cm, every node/.style={font=\small}]
\node[draw,rounded corners,minimum width=2.2cm] (detect) {Détecter};
\node[draw,rounded corners,minimum width=2.2cm,right=of detect] (rel) {Relier};
\node[draw,rounded corners,minimum width=2.2cm,right=of rel] (struct) {Structurer};
\node[draw,rounded corners,minimum width=2.2cm,right=of struct] (generate) {Générer};
\draw[-{Latex},thick] (detect)--(rel);
\draw[-{Latex},thick] (rel)--(struct);
\draw[-{Latex},thick] (struct)--(generate);
\end{tikzpicture}
\end{frame}

% ------------------------------------------------
\begin{frame}{Ce qui différencie réellement les approches}
\begin{itemize}
  \item \textbf{Qu'est-ce qui est détecté ?}
  \begin{itemize}
    \item éléments visuels ;
    \item régions ;
    \item groupes ;
    \item widgets ;
    \item ou simplement zones de l'image.
  \end{itemize}
  \item \textbf{Comment la relation parent--enfant est-elle déterminée ?}
  \begin{itemize}
    \item contenance géométrique ;
    \item découpage spatial ;
    \item raisonnement sémantique ;
    \item règles de mise en page ;
  \end{itemize}
  \item \textbf{Que doit représenter le modèle ?}
  \begin{itemize}
    \item la structure visuelle ;
    \item les composants ;
    \item la mise en page ;
    \item ou plusieurs de ces aspects.
  \end{itemize}
\end{itemize}
\end{frame}

% ------------------------------------------------
\begin{frame}{Lecture des méthodes par rapport à notre question}
\begin{block}{Notre objectif}
Obtenir une représentation intermédiaire suffisamment riche pour
permettre ensuite de générer du code frontend.
\end{block}

\vspace{0.2cm}
\begin{itemize}
  \item \method{Prototype2Code} montre l'intérêt d'un arbre orienté mise en page.
  \item \method{DCGen} montre qu'une hiérarchie peut émerger directement d'une décomposition de l'image.
  \item \method{ScreenCoder} montre l'intérêt de séparer \textbf{détection} et \textbf{construction de la structure}.
  \item \method{DesignCoder} montre l'intérêt d'ajouter la \textbf{sémantique} aux informations géométriques pour construire une hiérarchie de composants.
\end{itemize}

\begin{alertblock}{Question qui se dégage}
Quel niveau d'information doit contenir notre modèle pour être
indépendant de la technologie cible tout en restant suffisamment riche
pour générer ensuite React, HTML/CSS, Flutter, etc. ?
\end{alertblock}
\end{frame}

% ------------------------------------------------
\begin{frame}{Conclusion}
\begin{block}{Conclusion de l'étude}
Les travaux étudiés convergent vers une même idée :
\textbf{la maquette doit être transformée en une représentation structurée
avant de produire le code}.
\end{block}

\vspace{0.25cm}
Mais plusieurs stratégies existent :
\begin{itemize}
  \item \textbf{géométrique} : reconstruire les relations de contenance ;
  \item \textbf{spatiale} : découper récursivement l'image ;
  \item \textbf{structurelle} : convertir les régions en arbre de layout ;
  \item \textbf{sémantique} : comprendre les rôles et relations des composants.
\end{itemize}

\vspace{0.2cm}
\centering
\textbf{La problématique devient alors : quelle combinaison de ces informations
est nécessaire pour définir un modèle intermédiaire adapté à la génération
multi-technologies ?}
\end{frame}

% ------------------------------------------------
\begin{frame}{Références principales}
\footnotesize
\begin{thebibliography}{99}

\bibitem{prototype}
S. Xiao, Y. Chen, J. Li, L. Chen, T. Zhou.
\newblock \emph{Prototype2Code: End-to-end Front-end Code Generation from UI Design Prototypes}.
\newblock arXiv:2405.04975, 2024.
\newblock \url{https://arxiv.org/abs/2405.04975}

\bibitem{dcgen}
Y. Wan, C. Wang, Y. Dong, W. Wang, S. Li, Y. Huo, M. Lyu.
\newblock \emph{Divide-and-Conquer: Generating UI Code from Screenshots}.
\newblock Proc. ACM Software Engineering, FSE, 2025.
\newblock \url{https://doi.org/10.1145/3729364}



\bibitem{screencoder}
Y. Jiang, Y. Zheng, Y. Wan, J. Han, Q. Wang, M. R. Lyu, X. Yue.
\newblock \emph{ScreenCoder: Advancing Visual-to-Code Generation for Front-End Automation via Modular Multimodal Agents}.
\newblock arXiv:2507.22827, 2025.
\newblock \url{https://arxiv.org/abs/2507.22827}

\bibitem{designcoder}
Y. Chen, S. Ding, Y. Zhang, W. Chen, J. Du, L. Sun, L. Chen.
\newblock \emph{DesignCoder: Hierarchy-Aware and Self-Correcting UI Code Generation with Large Language Models}.
\newblock arXiv:2506.13663, 2025.
\newblock \url{https://arxiv.org/abs/2506.13663}

\end{thebibliography}
\end{frame}

\end{document}
